\subsection*{Texto A. Reseña académica}
\textit{Texto original de práctica. La obra y sus datos son ficticios.}

En \emph{Mapas del barrio}, de Elena Prado, la autora propone que los mapas escolares incorporen recorridos cotidianos y no solo avenidas principales. El libro presenta talleres realizados en dos grupos de secundaria de una misma ciudad. Antes del taller, los estudiantes dibujaron sus trayectos; después discutieron qué lugares habían omitido. En los dibujos finales aparecieron espacios que el plano oficial no señalaba: una biblioteca vecinal, un paso usado durante las lluvias y un punto de reunión.

La principal virtud del libro es mostrar que un mapa selecciona información según su propósito. No se afirma que cualquier representación sea adecuada para cualquier tarea: un itinerario personal puede servir para conversar sobre experiencias, pero no reemplaza un plano técnico para construir un puente. La explicación de los talleres es clara; en cambio, la autora ofrece pocos detalles sobre cómo eligió a los participantes. Por ello, los resultados permiten pensar otras maneras de enseñar cartografía, aunque no bastan para asegurar que todos los estudiantes aprenderían más con el mismo procedimiento.

\subsection*{Texto B. Informe de investigación}
\textit{Texto original de práctica; estudio y cifras ficticios para interpretación.}

Un equipo examinó si ofrecer un recordatorio modificaba la devolución puntual de libros. Participaron voluntariamente cuarenta personas de una biblioteca universitaria. Por sorteo, veinte recibieron un aviso dos días antes de la fecha límite y veinte no lo recibieron. Ambos grupos conservaron las mismas condiciones de préstamo. En el grupo con recordatorio, dieciséis devolvieron a tiempo; en el otro, doce lo hicieron.

El equipo describió una diferencia de veinte puntos porcentuales entre los grupos. Señaló que el resultado era compatible con una posible utilidad del recordatorio, pero evitó presentarlo como prueba definitiva. El número de participantes fue pequeño y todos procedían de una sola biblioteca. Además, no se registró si cada persona abrió el aviso. Se propuso repetir el estudio con más participantes y registrar la recepción efectiva del mensaje. El informe distinguió así entre observar una diferencia en esta muestra y demostrar que la misma intervención funcionará en cualquier biblioteca.
